\documentclass[reqno,11pt]{amsart}

\usepackage[margin=1in]{geometry}
\usepackage[T1]{fontenc}
\usepackage{lmodern}
\usepackage{microtype}
\usepackage{amsmath,amssymb,amsthm,mathtools}
\usepackage[dvipsnames]{xcolor}
\usepackage[most]{tcolorbox}
\usepackage{tikz}
\usepackage{enumitem}
\usepackage[hidelinks]{hyperref}
\hypersetup{
  pdftitle={Lecture 2, MATH 327: Examples and Subgroups},
  pdfauthor={Manish Patnaik},
  pdfsubject={MATH 327 course notes}
}

\definecolor{courseblue}{HTML}{1F4E79}
\definecolor{lightblue}{HTML}{EAF2F8}
\definecolor{darkgray}{HTML}{333333}
\definecolor{definitiongreen}{HTML}{EEF8EE}
\definecolor{definitionborder}{HTML}{8FB58F}
\definecolor{lightpink}{HTML}{FCE8EF}

\setlength{\parindent}{0pt}
\setlength{\parskip}{0.65em}
\setlist[itemize]{topsep=0.25em,itemsep=0.15em}

% Dotted leaders in the table of contents.
\makeatletter
\def\@tocline#1#2#3#4#5#6#7{\relax
  \ifnum #1>\c@tocdepth
  \else
    \par\addpenalty\@secpenalty\addvspace{#2}%
    \begingroup\hyphenpenalty\@M
    \@ifempty{#4}{%
      \@tempdima\csname r@tocindent\number#1\endcsname\relax
    }{%
      \@tempdima#4\relax
    }%
    \parindent\z@\leftskip#3\relax\advance\leftskip\@tempdima\relax
    \rightskip\@pnumwidth plus4em\parfillskip-\@pnumwidth
    #5\leavevmode\hskip-\@tempdima
    \ifcase #1
      \or\or\hskip 1em\or\hskip 2em\else\hskip 3em
    \fi
    #6\nobreak\relax
    \dotfill\hbox to\@pnumwidth{\@tocpagenum{#7}}\par
    \nobreak
    \endgroup
  \fi}
\makeatother

\theoremstyle{definition}
\newtheorem{definition}{Definition}[section]
\tcolorboxenvironment{definition}{
  enhanced,
  breakable,
  colback=definitiongreen,
  colframe=definitionborder,
  boxrule=0.5pt,
  arc=1mm,
  left=2mm,
  right=2mm,
  top=1mm,
  bottom=1mm,
  before skip=0.8em,
  after skip=0.8em
}
\newtheorem{example}{Example}[section]
\theoremstyle{plain}
\newtheorem{theorem}{Theorem}[section]
\newtheorem{lemma}{Lemma}[section]
\newtheorem{proposition}{Proposition}[section]
\theoremstyle{remark}
\newtheorem*{remark}{Remark}
\newtheorem{exercise}[theorem]{Exercise}

\tcolorboxenvironment{exercise}{
	enhanced,
	colback=lightpink,
	colframe=WildStrawberry,
	boxrule=0.7pt,
	arc=2pt,
	left=6pt,
	right=6pt,
	top=5pt,
	bottom=5pt
}


\numberwithin{equation}{section}

\newcommand{\Z}{\mathbb{Z}}
\newcommand{\R}{\mathbb{R}}
\newcommand{\GL}{\operatorname{GL}}
\newcommand{\id}{\operatorname{id}}
\newcommand{\be}{\begin{eqnarray}}
\newcommand{\ee}{\end{eqnarray}}

\begin{document}

\begin{center}
  \colorbox{lightblue}{%
    \begin{minipage}{0.94\textwidth}
      \vspace{0.55em}
      \begin{tabular*}{\textwidth}{@{\extracolsep{\fill}}lr}
        {\color{courseblue}\bfseries\LARGE Lecture 2} &
        {\color{courseblue}\bfseries\LARGE MATH 327}\\[0.35em]
         & \textbf{Date: Sep. 4, 2026}\\
        & \textbf{Instructor: Manish Patnaik}
      \end{tabular*}
      \vspace{0.45em}
    \end{minipage}}
\end{center}

\setcounter{tocdepth}{1}
\tableofcontents

\section{Review }

Recall that a group is a triple $(G,\cdot,e)$ consisting of a set $G$, a
composition law $\cdot$, and an identity element $e$. Depending on the
context, the identity may also be denoted by $1$, $e_G$, or $1_G$.

We also defined, for each $x\in G$, the cyclic subgroup $\langle x\rangle$,
which consists of all powers (positive, negative, and zero) of $x$. The
cardinality of $\langle x\rangle$ is also called the order of $x$ (and it can
be infinite).

Let us now turn to examples.


\section{Examples}

\begin{example}[The integers under addition]
The triple
\be
(\Z,+,0)
\ee
 is a group. Its composition law is
\be
\Z\times\Z\longrightarrow\Z,
  \qquad (m,n)\longmapsto m+n.
\ee
The identity element is $0$, and the inverse of $m$ is $-m$.
\end{example}

\begin{example}[The integers under multiplication: a non-example]
The triple
\be
(\Z,\times,1)
\ee
is \emph{not} a group, although it is a semigroup. Multiplication is
associative and $1$ is an identity, so axioms \textbf{(1)} and \textbf{(2)}
hold. Axiom \textbf{(3)} fails because most integers do not have
multiplicative inverses in $\Z$.
\end{example}


\begin{example}[Integers modulo $n$]
Let $n$ be a positive integer. The set of residue classes modulo $n$ is
\be
\Z/n\Z=\{\overline{0},\overline{1},\overline{2},\ldots,
  \overline{n-1}\}.
\ee
For example,
\begin{align*}
  \overline{0}&=\{0,\pm n,\pm 2n,\pm 3n,\ldots\},\\
  \overline{1}&=\{1,1\pm n,1\pm 2n,1\pm 3n,\ldots\}.
\end{align*}
Every integer can be written $m=an+r$ with $0\le r<n$, and then
$\overline{m}=\overline{r}$. Addition is defined by
\be
\overline{m}+\overline{\ell}=\overline{m+\ell}.
\ee
Thus
\be
(\Z/n\Z,+,\overline{0})
\ee
is a group, called the \emph{cyclic group of order $n$}.
\end{example}
%
%\begin{definition}[Cyclic group]
%A group $G$ is \emph{cyclic} if there is an element $x\in G$ such that
%\be
%G=\{x^k:k\in\Z\}.
%\ee
%In this case, $x$ is called a \emph{generator} of $G$.
%\end{definition}

In the preceding examples that are groups, the group operation is both
associative and commutative.


\begin{example}[The general linear group]
The set
\be
\GL_n(\R)=\{A:A\text{ is an invertible }n\times n
  \text{ matrix with entries in }\R\}
\ee
is a group under matrix multiplication. Its identity is the identity matrix
\be
I_n=
  \begin{pmatrix}
	    1 &        & 0\\
	      & \ddots &  \\
	    0 &        & 1
	  \end{pmatrix}.
\ee
Thus the group is written $(\GL_n(\R),\cdot,I_n)$, and the composition law
sends $(A,B)$ to $AB$. Note that the group $\GL_n(\R)$ carries additional
geometric structure: it is a topological space, a differentiable manifold,
and a real algebraic group. In particular, it is a \emph{Lie group}.
\end{example}

\newcommand{\SL}{\mathrm{SL}}

\begin{example}[The special linear group]
Consider the following subset of the previous example:
\be
\SL_n(\R)=\{A\in M_n(\R):\det(A)=1\}.
\ee
It is a group under matrix multiplication. Indeed, $I_n\in\SL_n(\R)$, and
if $A,B\in\SL_n(\R)$, then
\be
\det(AB)=\det(A)\det(B)=1,
\ee
so $AB\in\SL_n(\R)$. Moreover,
\be
\det(A^{-1})=\det(A)^{-1}=1,
\ee
so $A^{-1}\in\SL_n(\R)$. Associativity is inherited from matrix
multiplication. Thus $(\SL_n(\R),\cdot,I_n)$ is a group.
\end{example}



\section{Subgroups}

\begin{definition}[Subgroup]
Let $(G,\cdot,e)$ be a group. A subset $H\subseteq G$ is called a
\emph{subgroup} of $G$ if $H$, under the restriction of the operation on
$G$, is itself a group.
\end{definition}

Not all subsets are groups. For example, the subset $\{0,3\}\subset\Z$ is
not a group under addition, since it is not closed under addition.

\begin{proposition}
Let $(G,\cdot,e)$ be a group, and let $H\subseteq G$ be a subset satisfying
the following conditions:
\begin{enumerate}
  \item $e\in H$.
  \item $H$ is closed under the group operation: if $h_1,h_2\in H$, then
  \be
  h_1\cdot h_2\in H.
  \ee
  \item $H$ is closed under inverses: if $h\in H$, then
  \be
  h^{-1}\in H.
  \ee
\end{enumerate}
Then $H$ is a subgroup of $G$.
\end{proposition}

\begin{proof}
The operation on $H$ is the restriction of the operation on $G$.
Associativity is therefore inherited from $G$, so it does not need to be
checked separately. By condition (2) $H$ is closed under this operation, and
since $e$ is an identity in $G$, it also satisfies the identity property of
$H$. Finally, condition (3) ensures that the inverse in $G$ of every element
of $H$ also belongs to $H$.
\end{proof}

\begin{remark}
Suppose $H\subseteq G$ is finite. If $e\in H$ and $H$ is closed under the
group operation, then $H$ is automatically closed under inverses. To see
this, fix $h\in H$ and consider the map
\be
L_h\colon H\longrightarrow H,
  \qquad x\longmapsto hx.
\ee
This map is well-defined because $H$ is closed under the group operation. It
is injective: if $hx=hy$, then multiplying on the left by $h^{-1}$ in $G$
gives $x=y$. Since $H$ is finite, $L_h$ is therefore surjective. Because
$e\in H$, there exists $y\in H$ such that $hy=e$. By uniqueness of inverses
in $G$, we have $y=h^{-1}$. Hence $h^{-1}\in H$.

\begin{exercise}This shows that for every $h \in H$ there exists $g$ such that $hg=e$, so $g$ is a right inverse. Why is $g$ also a left inverse? \end{exercise}
\end{remark}

\section{Permutation Groups}

Next we turn to a very important class of finite groups, the so-called permutation groups. You have probably seen these before, but I want to reintroduce the notation. Let 
\be
[n]:=\{1,2,\ldots,n\}.
\ee
A permutation of $[n]$ is a bijection $p\colon[n]\to[n]$. It may be displayed
in two-line notation as
\be
p=
\begin{pmatrix}
	1&2&\cdots&n\\
	p(1)&p(2)&\cdots&p(n)
\end{pmatrix}.
\ee

If $p$ and $q$ are permutations, we use the convention
\be
qp=q\circ p;
\ee
thus the rightmost permutation is applied first. Let us also note that if $p$
is a permutation, we write $p^{-1}$ for the inverse permutation; that is, if
$p(i)=k$, then $p^{-1}(k)=i$. One can check that
\be
pp^{-1}=p^{-1}p=\id,
\ee
where $\id$ is the identity permutation.

Let $S_n$ denote the set of all permutations of $[n]$.

\begin{proposition}
	The triple $(S_n,\circ,\id)$ is a group.
\end{proposition}

\begin{proof}
	The composite of two bijections is again a bijection, so $S_n$ is closed
	under composition. Composition of functions is associative, the identity
	map $\id$ is the identity element, and the inverse of a bijection is again a
	bijection. Hence all the group axioms hold.
\end{proof}



\end{document}
